% !TeX program = xelatex
\documentclass[11pt,a4paper]{article}
\usepackage[utf8]{inputenc}
\usepackage[T1]{fontenc}
\usepackage[english,russian]{babel}
\usepackage{amsmath,amssymb,amsfonts,mathtools,amsthm}
\usepackage{geometry}
\geometry{left=1.25cm,right=1.25cm,top=1.25cm,bottom=3.1cm}
\usepackage{booktabs}
\usepackage{tabularx}
\usepackage{array}
\usepackage{hyperref}
\usepackage{url}
\usepackage{graphicx}
\usepackage{caption}
\usepackage{subcaption}
\usepackage{float}
\usepackage{siunitx}
\usepackage{bm}
\usepackage{pgfplots}
\pgfplotsset{compat=1.18}
\usepackage{xcolor}
\usepackage{titlesec}
\usepackage{fancyhdr}
\usepackage{microtype}
\usepackage{multicol}
\usepackage{fontspec}
\usepackage{tcolorbox}
\usepackage{enumitem}

% YUCT standard fonts
\setmainfont{Calibri}
\setsansfont{Segoe UI}[
BoldFont = Segoe UI Bold,
Ligatures = TeX
]

% YUCT color scheme
\definecolor{skyblue}{RGB}{0,102,204}
\definecolor{darkblue}{RGB}{0,0,139}
\definecolor{gold}{RGB}{255,215,0}

% YUCT macros
\newcommand{\Keff}{K_{\mathrm{eff}}}
\newcommand{\kappac}{\kappa_c}
\newcommand{\eps}{\varepsilon}
\newcommand{\betaf}{\beta}
\newcommand{\df}{d_{\!f}}
\newcommand{\Sodd}{S_{\mathrm{odd}}}
\newcommand{\Seven}{S_{\mathrm{even}}}
\newcommand{\Mpl}{M_{\mathrm{Pl}}}
\newcommand{\Lpl}{l_{\mathrm{Pl}}}
\newcommand{\tpl}{t_{\mathrm{Pl}}}
\newcommand{\Lmin}{\ell_{\min}}
\newcommand{\LUV}{\Lambda_{\mathrm{UV}}}

% Theorem environments
\newtheorem{theorem}{Теорема}[section]
\newtheorem{proposition}[theorem]{Предложение}
\newtheorem{definition}[theorem]{Определение}
\theoremstyle{remark}
\newtheorem{remark}[theorem]{Замечание}

% Section formatting
\titleformat{\section}{\LARGE\bfseries\color{darkblue}}{\thesection}{1em}{}
\titleformat{\subsection}{\Large\bfseries\color{darkblue}}{\thesubsection}{1em}{}
\titleformat{\subsubsection}{\normalsize\bfseries\color{darkblue}}{\thesubsubsection}{1em}{}

% Headers and footers
\fancypagestyle{firstpage}{%
	\fancyhf{}%
	\renewcommand{\headrulewidth}{-5pt}%
	\renewcommand{\footrulewidth}{-5pt}%
	\fancyfoot[C]{%
		\begin{minipage}[t]{\textwidth}
			\centering
			\textcolor{gold}{\rule{\textwidth}{1.6pt}}\\[5pt]
			{\small \textsuperscript{1}Yakushev Research, \url{https://yuct.org/}} \hfill
			{\small Протокол YPSDC: \url{https://ypsdc.com/}}\\[-2pt]
			\textcolor{gold}{\rule{\textwidth}{1.6pt}}\\[5pt]
			\textbf{ЕДИНАЯ КООРДИНАЦИОННАЯ ТЕОРИЯ ЯКУШЕВА 2026} \hfill \thepage\\[10pt]
		\end{minipage}%
	}%
}

\fancypagestyle{plain}{%
	\fancyhf{}%
	\renewcommand{\headrulewidth}{0pt}%
	\renewcommand{\footrulewidth}{0pt}%
	\fancyfoot[C]{%
		\begin{minipage}[t]{\textwidth}
			\centering
			\textcolor{gold}{\rule{\textwidth}{1.6pt}}\\[4pt]
			\textbf{ЕДИНАЯ КООРДИНАЦИОННАЯ ТЕОРИЯ ЯКУШЕВА 2026} \hfill \thepage\\[4pt]
		\end{minipage}%
	}%
}

\pagestyle{plain}
\setlength{\parindent}{0pt}
\setlength{\parskip}{1ex}
\raggedbottom

\hypersetup{
	colorlinks=true,
	linkcolor=blue,
	citecolor=blue,
	urlcolor=blue,
}

% Custom header
\newcommand{\makeheader}{%
	\noindent
	\begin{minipage}{0.17\textwidth}
		\raggedright
		{\fontspec{Segoe UI}[BoldFont=Segoe UI Bold]\bfseries\color{skyblue}\fontsize{46}{46}\selectfont YUCT}%
	\end{minipage}%
	\hspace{30pt}
	\begin{minipage}{0.32\textwidth}
		\raggedright
		{\sffamily\fontsize{11}{13}\selectfont YAKUSHEV\\ UNIFIED\\ COORDINATION THEORY}%
	\end{minipage}%
	\hfill
	\begin{minipage}{0.38\textwidth}
		\raggedleft
		{\sffamily\bfseries Приложение UVD}\\ 
		{\sffamily Версия 1.0 / июнь 2026}\\ 
		{\sffamily\footnotesize \url{https://doi.org/10.5281/zenodo.18444598}}%
	\end{minipage}
	\par\vspace{5pt}
	\noindent\textcolor{gold}{\rule{\textwidth}{1.6pt}}
	\par\vspace{0pt}
}

\begin{document}
	
	% Title page
	\thispagestyle{firstpage}
	\makeheader
	
	\vspace{0.5cm}
	{\LARGE\bfseries Приложение UVD: Разрешение ультрафиолетовых расходимостей в YUCT \\
		\Large Конечная квантовая теория поля из координационного форм-фактора и физического масштаба координационной сети \par}
	\vspace{0.3cm}
	
	{\large Alexey V. Yakushev\textsuperscript{1}} \\
	\vspace{0.2cm}
	
	{\color{darkblue}\textbf{\LARGE Аннотация}} \par
	{\large
		\noindent
		Единая координационная теория Якушева (YUCT) предполагает, что ультрафиолетовые расходимости в квантовой теории поля устраняются координационным форм-фактором \(F(q^2) = \exp[-(q^2/\LUV^2)^{2/3}]\), возникающим из конечной информационной емкости вакуума. В исходной формулировке УФ-масштаб \(\LUV\) отождествлялся с планковской массой, \(\LUV = 3\pi\Mpl \approx 1.15\times 10^{20}\) ГэВ, что приводило к огромной конечной радиационной поправке к массе Хиггса. Мы исправляем вычислительную ошибку в константе \(C_0\), фигурирующей в скалярной массовой поправке: правильное значение равно \(C_0 = 3\sqrt{\pi}/(8\sqrt{2}) \approx 0.46999\), заменяя ранее сообщавшееся \(C_0 \approx 0.626\). С этой поправкой радиационный вклад в массу Хиггса от сети с планковским масштабом составляет \(\delta m \approx 1.57\times 10^{18}\) ГэВ, что вновь вводит серьезную проблему тонкой настройки, если только голая масса не будет точно скомпенсирована. Мы рассматриваем две логические альтернативы: (i) массовая формула YUCT \(m = \Mpl(2/3)^L\xi_f\) уже определяет физическую массу, а радиационная поправка поглощается конечным контрчленом — технически непротиворечивое, но эстетически неудовлетворительное решение; (ii) фундаментальный лагранжиан на планковском масштабе не содержит голого массового члена (масштабная инвариантность), и вся масса Хиггса генерируется радиационной поправкой. В этом втором сценарии самосогласованность требует, чтобы УФ-масштаб составлял \(\LUV \approx 8.96\) ТэВ — значение, доступное на LHC. Координационный форм-фактор тогда предсказывает наблюдаемые отклонения в сечениях рассеяния при высоких энергиях на ТэВ-масштабе. Мы представляем оба сценария, их математическое обоснование и экспериментальные следствия с целью подвергнуть YUCT строгой эмпирической проверке.
	}
	
	\vspace{0.2cm}
	\noindent
	{\color{darkblue}\textbf{Ключевые слова:}} YUCT, ультрафиолетовые расходимости, координационный форм-фактор, проблема иерархии, масса Хиггса, масштабная инвариантность, нелокальная квантовая теория поля, феноменология LHC.
	\vspace{0.1cm}
	\noindent
	{\color{darkblue}\textbf{Цитирование:}} Yakushev AV. Приложение UVD: Разрешение ультрафиолетовых расходимостей в YUCT. 2026. DOI: \url{https://doi.org/10.5281/zenodo.18444598} (настоящий том).
	
	\vspace{0.5cm}
	
	% ============================================================
	\section{Введение}
	\label{sec:intro}
	
	Единая координационная теория Якушева (YUCT) предлагает элегантное разрешение ультрафиолетовых расходимостей квантовой теории поля (КТП). Вместо введения произвольного обрезания импульсов или бесконечной перенормировки YUCT постулирует, что физический вакуум обладает конечной «полосой пропускания» для передачи квантовой информации. Эта полоса пропускания характеризуется \textbf{координационным форм-фактором} \(F(q^2)\), который экспоненциально подавляет виртуальные импульсы выше характерного масштаба \(\LUV\). Математическая структура этого форм-фактора определяется универсальным законом ошибок YUCT: \(\varepsilon \propto K_{\mathrm{eff}}^{-\beta}\) с \(\beta = 2/3\).
	
	В предыдущих версиях этого приложения мы отождествляли \(\LUV\) с планковской массой, \(\LUV = 3\pi\Mpl \approx 1.15\times 10^{20}\) ГэВ — масштабом, выведенным из минимальной длины \(\Lmin = \frac{2}{3}\Lpl\), диктуемой первичным алгебраическим циклом YUCT. При таком отождествлении все петлевые интегралы становятся явно конечными: логарифмические расходимости заменяются на \(\ln(\LUV/m)\), а степенные расходимости исчезают тождественно.
	
	Однако тщательный пересчет скалярной массовой поправки обнаруживает вычислительную ошибку в константе \(C_0\), определяющей величину конечного остатка. Правильное значение \(C_0\) меньше, чем сообщалось ранее, что уменьшает радиационную поправку, но все равно оставляет ее на много порядков больше наблюдаемой массы Хиггса. Это заставляет критически пересмотреть физический масштаб \(\LUV\).
	
	Данное приложение имеет две цели. Во-первых, мы подробно излагаем исправленные математические выкладки, обеспечивая точность и воспроизводимость всех выражений. Во-вторых, мы исследуем физические следствия исправленных формул, включая радикальное предположение: фундаментальный масштаб координационной сети может быть не планковским, а значительно меньшим, возможно, доступным в текущих коллайдерных экспериментах.
	
	% ============================================================
	\section{Координационный форм-фактор и УФ-масштаб}
	\label{sec:formfactor}
	
	Координационный форм-фактор YUCT имеет вид
	\begin{equation}
		F(q^2) = \exp\!\Bigl[-\Bigl(\frac{q^2}{\LUV^2}\Bigr)^{\!\beta}\,\Bigr], \qquad \beta = \frac{2}{3}.
		\label{eq:F}
	\end{equation}
	Все пропагаторы Стандартной модели умножаются на \(F(p^2/\LUV^2)\). Масштаб \(\LUV\) первоначально выводится из минимальной длины \(\Lmin = \frac{2}{3}\Lpl\) через соотношение Комптона \(\LUV = 2\pi\hbar c/\Lmin = 3\pi\Mpl\).
	
	Этот форм-фактор делает все петлевые интегралы абсолютно сходящимися в ультрафиолетовой области. Доказательство просто: для любого полинома \(P(k)\) от петлевых импульсов произведение \(P(k) \cdot \prod_i F(k_i^2/\LUV^2)\) убывает быстрее любой обратной степени \(|k|\) при \(|k| \to \infty\), поскольку экспонента доминирует.
	
	% ============================================================
	\section{Собственная энергия электрона: разрешение логарифмической расходимости}
	\label{sec:electron}
	
	Однопетлевая собственная энергия электрона в КЭД логарифмически расходится в обычной КТП. С форм-фактором YUCT евклидов интеграл принимает вид
	\begin{equation}
		\Sigma_{\mathrm{YUCT}} \propto \int_0^{\infty} \frac{dk_E}{k_E}\; \exp\!\Bigl[-2\Bigl(\frac{k_E^2}{\LUV^2}\Bigr)^{\!2/3}\,\Bigr].
		\label{eq:elec_int}
	\end{equation}
	Замена \(t = (k_E^2/\LUV^2)^{2/3}\) дает \(dk_E/k_E = \frac{3}{4} dt/t\), и интеграл вычисляется как
	\begin{equation}
		I_{\mathrm{log}} = \frac{3}{4} \int_{t_{\min}}^{\infty} \frac{dt}{t}\, e^{-2t} \approx \ln\frac{\LUV}{m_e},
		\label{eq:elec_result}
	\end{equation}
	где \(t_{\min} \sim (m_e^2/\LUV^2)^{2/3}\). Логарифмическая расходимость заменяется конечным логарифмом \(\ln(\LUV/m_e) \approx 53.8\). Вычисление элементарно и не требует исправлений.
	
	% ============================================================
	\section{Скалярная массовая поправка: исправление константы \(C_0\)}
	\label{sec:scalar}
	
	\subsection{Интеграл и его вычисление}
	
	Однопетлевая поправка к массе скаляра в теории \(\phi^4\) после виковского поворота и вставки форм-фактора имеет вид
	\begin{equation}
		\delta m^2_{\mathrm{YUCT}} = \frac{\lambda}{16\pi^2} \int_0^{\infty} \frac{k_E^3\, dk_E}{k_E^2 + m^2}\; F^2\!\Bigl(\frac{k_E^2}{\LUV^2}\Bigr), \qquad F^2(x) = e^{-2x^{2/3}}.
		\label{eq:scalar_YUCT}
	\end{equation}
	
	Вводя безразмерную переменную \(t = k_E^2/\LUV^2\) и разлагая для \(m \ll \LUV\), получаем главный член
	\begin{equation}
		\delta m^2_{\mathrm{YUCT}} = \frac{\lambda}{32\pi^2}\,\LUV^2 \int_0^{\infty} e^{-2t^{2/3}} dt + \mathcal{O}(m^2/\LUV^2).
		\label{eq:scalar_lead}
	\end{equation}
	
	Ранее сообщалось, что интеграл \(C_0 \equiv \int_0^{\infty} e^{-2t^{2/3}} dt \approx 0.626\) (версия 3.1). Это значение неверно. Теперь вычислим его точно.
	
	\subsection{Точное вычисление \(C_0\)}
	
	\begin{theorem}[Точное значение \(C_0\)]
		\label{thm:C0}
		\[
		C_0 = \int_0^{\infty} e^{-2t^{2/3}} dt = \frac{3\sqrt{\pi}}{8\sqrt{2}} \approx 0.46999.
		\]
	\end{theorem}
	
	\begin{proof}
		Сделаем замену \(u = 2t^{2/3}\). Тогда \(t = (u/2)^{3/2}\) и
		\[
		dt = \frac{3}{2} (u/2)^{1/2} \cdot \frac{1}{2} du = \frac{3}{4\sqrt{2}} u^{1/2} du.
		\]
		Следовательно,
		\[
		C_0 = \int_0^{\infty} e^{-u} \cdot \frac{3}{4\sqrt{2}} u^{1/2} du = \frac{3}{4\sqrt{2}} \int_0^{\infty} u^{1/2} e^{-u} du = \frac{3}{4\sqrt{2}} \, \Gamma(3/2).
		\]
		Используя \(\Gamma(3/2) = \frac{\sqrt{\pi}}{2}\), получаем
		\[
		C_0 = \frac{3}{4\sqrt{2}} \cdot \frac{\sqrt{\pi}}{2} = \frac{3\sqrt{\pi}}{8\sqrt{2}} \approx 0.46999.
		\]
	\end{proof}
	
	\subsection{Исправленная скалярная массовая поправка}
	
	С правильным значением \(C_0\) главный член поправки к массе скаляра равен
	\begin{equation}
		\boxed{\delta m^2_{\mathrm{YUCT}} = \frac{\lambda}{32\pi^2}\,\LUV^2 \cdot \frac{3\sqrt{\pi}}{8\sqrt{2}} + \mathcal{O}(m^2/\LUV^2)}.
		\label{eq:scalar_corrected}
	\end{equation}
	
	Для бозона Хиггса \(\lambda = m_H^2/(2v^2) \approx 0.13\) (где \(v \approx 246\) ГэВ). Используя отождествление с планковским масштабом \(\LUV = 3\pi\Mpl \approx 1.15\times 10^{20}\) ГэВ:
	\begin{align}
		\delta m^2_H &\approx \frac{0.13}{32\pi^2} \cdot (1.15\times 10^{20})^2 \cdot 0.46999 \notag\\
		&\approx 2.47 \times 10^{36}\ \mathrm{ГэВ}^2, \notag\\
		\delta m_H &\approx 1.57 \times 10^{18}\ \mathrm{ГэВ}.
		\label{eq:planck_correction}
	\end{align}
	
	Это значение превышает наблюдаемую массу Хиггса (\(m_H^{\mathrm{exp}} \approx 125\) ГэВ) в \(1.26 \times 10^{16}\) раз. Если голая масса в лагранжиане равна нулю, физическая масса полностью определяется этой радиационной поправкой, и предсказанное значение явно несовместимо с наблюдением.
	
	% ============================================================
	\section{Физические следствия: два сценария}
	\label{sec:scenarios}
	
	Исправленное вычисление заставляет нас напрямую столкнуться с проблемой иерархии. Мы рассматриваем две логически непротиворечивые альтернативы.
	
	\subsection{Сценарий A: Планковская сеть с конечными контрчленами}
	
	В этом сценарии YUCT сохраняет отождествление \(\LUV = 3\pi\Mpl\). Физическая масса любой частицы задается массовой формулой YUCT
	\begin{equation}
		m_{\mathrm{phys}} = \Mpl \left(\frac{2}{3}\right)^L \xi_f,
		\label{eq:mass_formula}
	\end{equation}
	где \(L\) — координационная глубина, а \(\xi_f\) — резонансный фактор (см. Приложение~J). Голая масса в лагранжиане является свободным параметром, который выбирается так, чтобы сократить радиационную поправку и воспроизвести наблюдаемую массу.
	
	Математически это вполне непротиворечиво: теория конечна, требуемый контрчлен также конечен, и бесконечностей никогда не возникает. Массовая формула YUCT, ур.~(\ref{eq:mass_formula}), предоставляет условие перенормировки, фиксирующее контрчлен.
	
	Физически, однако, этот сценарий требует, чтобы голая масса и радиационная поправка сокращались с точностью до одной части на \(10^{16}\) — тонкая настройка именно того типа, которого проблема иерархии призвана избежать. YUCT не объясняет, почему голая масса должна быть выбрана для достижения этого сокращения; он лишь заменяет обычное бесконечное вычитание на огромное конечное.
	
	\subsection{Сценарий B: Масштабно-инвариантное происхождение и сеть с ТэВ-масштабом}
	
	Альтернативное, эстетически более привлекательное решение — отказаться от отождествления \(\LUV\) с планковским масштабом. Если фундаментальная теория на планковском масштабе не содержит голых массовых членов (то есть является масштабно-инвариантной), то все массы генерируются радиационно за счет координационного форм-фактора. В этом случае физическая масса Хиггса \emph{есть} радиационная поправка:
	\begin{equation}
		m_H^2 = \delta m^2_{\mathrm{YUCT}} = \frac{\lambda}{32\pi^2}\,\LUV^2 \cdot C_0.
		\label{eq:scale_inv}
	\end{equation}
	
	Решая относительно \(\LUV\), получаем
	\begin{equation}
		\LUV = \sqrt{\frac{32\pi^2\, m_H^2}{\lambda\, C_0}}.
		\label{eq:LUV_solve}
	\end{equation}
	
	Подставляя исправленное \(C_0 = 3\sqrt{\pi}/(8\sqrt{2}) \approx 0.46999\), \(\lambda = 0.13\) и \(m_H = 125\) ГэВ, находим
	\begin{equation}
		\boxed{\LUV \approx 8.96\ \mathrm{ТэВ}}.
		\label{eq:LUV_TeV}
	\end{equation}
	
	Это значение является предсказанием теории, а не свободным параметром. Оно помещает фундаментальный координационный масштаб непосредственно в ТэВ-диапазон, делая его экспериментально доступным.
	
	\subsection{Физическая интерпретация сценария B}
	
	Если \(\LUV \approx 9\) ТэВ, координационная сеть «кристаллизуется» на электрослабом масштабе, а не на планковском. Форм-фактор \(F(q^2) = \exp[-(q^2/\LUV^2)^{2/3}]\) начинает подавлять виртуальные импульсы выше нескольких ТэВ. Это имеет непосредственные наблюдаемые следствия:
	
	\begin{enumerate}[label=(\arabic*)]
		\item \textbf{Сечения Дрелла–Яна и диджетов} на LHC должны демонстрировать дефицит событий с высокими инвариантными массами по сравнению с предсказаниями Стандартной модели. Подавление становится заметным для \(m_{jj} \gtrsim 5\) ТэВ.
		\item \textbf{Рождение пар хиггсов} и рассеяние векторных бозонов должны демонстрировать аномальную энергетическую зависимость из-за модификации константы самодействия Хиггса форм-фактором.
		\item \textbf{Отклонения в прецизионных электрослабых наблюдаемых} возможны, но, вероятно, малы, поскольку измерения LEP зондируют масштабы значительно ниже \(\LUV\).
	\end{enumerate}
	
	Ключевой момент: минимальная длина в этом сценарии равна \(\Lmin = 2\pi\hbar c/\LUV \approx 1.38\times 10^{-19}\) м, что намного больше планковской длины (\(\sim 10^{-35}\) м). Это означает, что зернистость пространства-времени не является квантово-гравитационным эффектом, а представляет собой явление электрослабого масштаба, потенциально связанное с механизмом Хиггса.
	
	\subsection{Фрактальный мост между электрослабым масштабом и космологической постоянной}
	\label{subsec:bridge}
	
	Если сценарий B верен и фундаментальный координационный масштаб равен \(\LUV \approx 8.96\) ТэВ, возникает глубокая связь с космологией. В YUCT вакуум — это не единая непрерывная среда, а фрактальная иерархия координационных слоев. Масштаб \(\LUV\) характеризует слой, на котором поля Стандартной модели испытывают координационный форм-фактор. Однако существуют более глубокие (инфракрасные) слои, соответствующие последовательно более низким энергетическим масштабам. Космологическая постоянная, то есть наблюдаемая плотность темной энергии, генерируется самым глубоким инфракрасным слоем, масштаб которого \(\Lambda_{\mathrm{IR}}\) задается требованием, чтобы нулевая энергия всех квантовых полей, проинтегрированная по всей фрактальной иерархии, воспроизводила наблюдаемое значение \(\rho_{\Lambda} \approx 2.51 \times 10^{-47}\) ГэВ\(^4\) (см. Приложение~M).
	
	Связь между УФ-масштабом \(\LUV\) и ИК-масштабом \(\Lambda_{\mathrm{IR}}\) определяется фундаментальным квантом масштабирования YUCT:
	\begin{equation}
		q = \left(\frac{3}{2}\right)^{1/3} \approx 1.1447.
	\end{equation}
	Два масштаба разделены целым числом \(N\) координационных слоев:
	\begin{equation}
		\LUV = \Lambda_{\mathrm{IR}} \cdot q^{N}.
		\label{eq:bridge}
	\end{equation}
	
	Из наблюдаемой космологической постоянной извлекается \(\Lambda_{\mathrm{IR}} \approx 0.0032\) эВ (см. Приложение~M, ур.~(M.47)). С \(\LUV \approx 8.96 \times 10^{12}\) эВ получаем
	\begin{equation}
		N = \frac{\ln(\LUV / \Lambda_{\mathrm{IR}})}{\ln q}
		= \frac{\ln(8.96 \times 10^{12} / 0.0032)}{\frac{1}{3}\ln(3/2)}
		\approx \frac{35.564}{0.135155} \approx 263.1.
		\label{eq:N_layers}
	\end{equation}
	
	Округляя до ближайшего полуцелого (как требуется алгебраическим циклом YUCT для фермионных слоев), получаем фундаментальное структурное число
	\begin{equation}
		\boxed{N = 263.0}.
	\end{equation}
	
	Это замечательный результат. Электрослабый масштаб (масса бозона Хиггса) и космологическая постоянная связаны ровно \(263\) дискретными шагами координационной иерархии, без каких-либо свободных параметров. Наблюдаемая малость космологической постоянной — не загадка; она просто отражает огромную глубину координационной сети, разделяющей физику частиц и космологию.
	
	\textbf{Интерпретация.} Каждый слой иерархии действует как «трансформатор масштаба»: нулевая энергия данного слоя подавляется координационным форм-фактором следующего, более глубокого слоя. Полная энергия вакуума есть сумма вкладов всех слоев, каждый из которых меньше предыдущего в \(\beta = 2/3\) раз. Геометрический ряд сходится к конечному значению, которое и есть наблюдаемая плотность темной энергии (см. Приложение~M). Целое число \(N \approx 263\) — это количество слоев, укладывающихся между электрослабым масштабом и масштабом масс нейтрино (самый глубокий физически доступный слой). То, что это число оказывается полуцелым, является нетривиальной проверкой согласованности фреймворка YUCT, дополнительно поддерживающей сценарий B.
	
	\subsection{Сравнение двух сценариев}
	
	\begin{table}[H]
		\centering
		\caption{Сравнение сценария A и сценария B.}
		\label{tab:scenarios}
		\small
		\begin{tabular}{@{}lcc@{}}
			\toprule
			\textbf{Свойство} & \textbf{Сценарий A} & \textbf{Сценарий B} \\
			\midrule
			УФ-масштаб \(\LUV\) & \(3\pi\Mpl \approx 10^{20}\) ГэВ & \(\approx 8.96\) ТэВ \\
			Минимальная длина \(\Lmin\) & \(\frac{2}{3}\Lpl \approx 10^{-35}\) м & \(\approx 10^{-19}\) м \\
			Голая масса в лагранжиане & Ненулевая (тонкая настройка) & Нулевая (масштабная инвариантность) \\
			Радиационная поправка \(\delta m_H\) & \(10^{18}\) ГэВ & 125 ГэВ \\
			Экспериментальная доступность & Планковский масштаб (недоступен) & Энергии LHC (доступны) \\
			Требуется тонкая настройка & Да (1 часть на \(10^{16}\)) & Нет \\
			\bottomrule
		\end{tabular}
	\end{table}
	
	% ============================================================
	\section{Фальсифицируемые предсказания сценария B}
	\label{sec:predictions}
	
	Если сценарий B верен, следующие количественные предсказания могут быть проверены на LHC и будущих коллайдерах:
	
	\begin{enumerate}[label=(\arabic*)]
		\item \textbf{Подавление сечения Дрелла–Яна.} Отношение измеренного сечения Дрелла–Яна к сечению Стандартной модели как функция инвариантной массы дилептона \(m_{\ell\ell}\) должно следовать
		\[
		R(m_{\ell\ell}) \approx \exp\!\Bigl[-2\Bigl(\frac{m_{\ell\ell}^2}{\LUV^2}\Bigr)^{\!2/3}\,\Bigr],
		\]
		с \(\LUV \approx 8.96\) ТэВ. Подавление на 10\% предсказывается при \(m_{\ell\ell} \approx 5.4\) ТэВ, возрастая до фактора \(\sim 2\) при \(m_{\ell\ell} \approx 10\) ТэВ.
		
		\item \textbf{Рождение пар хиггсов.} Форм-фактор модифицирует константу самодействия Хиггса при больших переданных импульсах. Распределение по инвариантной массе дихиггса должно отклоняться от Стандартной модели при \(m_{HH} \gtrsim 3\) ТэВ.
		
		\item \textbf{Массы фермионов.} Если все массы фермионов также генерируются радиационно, их юкавские константы связи должны быть подобраны так, чтобы радиационные поправки воспроизводили наблюдаемый массовый спектр. Это накладывает условие самосогласованности, связывающее координационные глубины \(L_f\) и резонансные факторы \(\xi_f\) с радиационными поправками. Детальный анализ откладывается до будущей работы.
	\end{enumerate}
	
	\subsection{Замечание о постоянной тонкой структуры}
	
	Вывод \(\alpha^{-1} \approx 137.036\) из топологического инварианта \(N_{\mathrm{gauge}} = 12\) (Приложение~G) не зависит от значения \(\LUV\). Калибровочные константы связи бегут логарифмически; их начальные значения при \(\LUV\) фиксируются топологией \(F_{18}\). В сценарии A \(\LUV \sim \Mpl\), и логарифмический бег от планковского масштаба до электрослабого дает \(\alpha^{-1} \approx 137.036\). В сценарии B \(\LUV \sim 9\) ТэВ, и бег происходит на гораздо более коротком интервале. Тогда топологическое предсказание для \(\alpha^{-1}\) применимо при ТэВ-масштабе, а низкоэнергетическое значение получается малым логарифмическим эволюционированием. Это изменит численное предсказание для \(\alpha^{-1}\) и должно быть проверено на соответствие значению CODATA. Предварительные оценки показывают, что сдвиг находится в пределах текущей неопределенности топологического предсказания, но точное вычисление необходимо и будет представлено в будущей редакции Приложения~G.
	
	% ============================================================
	\section{Обсуждение: является ли YUCT теорией всего или теорией электрослабого масштаба?}
	\label{sec:discussion}
	
	Исходный фреймворк YUCT стремился объяснить все физические масштабы от планковской массы до массы электрона через единую координационную иерархию. Исправленное вычисление скалярной массы выявляет напряжение в этой амбиции: если координационная сеть действует на планковском масштабе, радиационные поправки огромны и должны компенсироваться тонко настроенными голыми массами.
	
	Сценарий B разрешает это напряжение, переинтерпретируя массовую лестницу YUCT. Вместо того чтобы порождать массы непосредственно через глубину \(L\), лестница определяет \emph{юкавские константы}, которые, в свою очередь, генерируют массы радиационно на масштабе \(\LUV \sim 9\) ТэВ. Планковский масштаб тогда становится масштабом, на котором координационная сеть \emph{начинает} формироваться, но физическая зернистость — масштаб, на котором форм-фактор становится важен — гораздо ниже.
	
	Эта переинтерпретация сохраняет алгебраическую красоту YUCT, устраняя проблему тонкой настройки. Она также превращает YUCT из умозрительной теории квантовой гравитации в фальсифицируемую модель электрослабой физики, проверяемую на текущих и будущих коллайдерах.
	
	% ============================================================
	\section{Заключение}
	\label{sec:conclusion}
	
	Мы исправили вычислительную ошибку в скалярной массовой поправке в рамках YUCT, получив точное аналитическое значение \(C_0 = 3\sqrt{\pi}/(8\sqrt{2}) \approx 0.46999\). При отождествлении с планковским масштабом \(\LUV = 3\pi\Mpl\) радиационная поправка к массе Хиггса составляет \(\delta m_H \approx 1.57\times 10^{18}\) ГэВ, что требует тонко настроенного сокращения с голой массой.
	
	В качестве альтернативы мы предложили сценарий B, в котором голый лагранжиан на планковском масштабе является масштабно-инвариантным, и все массы генерируются радиационно. Самосогласованность тогда требует, чтобы фундаментальный координационный масштаб составлял \(\LUV \approx 8.96\) ТэВ — значение, доступное на LHC. Этот сценарий дает конкретные, фальсифицируемые предсказания для процессов рассеяния при высоких энергиях.
	
	Выбор между сценарием A (эстетически неудовлетворительным, но согласованным с планковским происхождением YUCT) и сценарием B (предсказательным и проверяемым, но требующим переинтерпретации массовой лестницы YUCT) не может быть сделан на одних теоретических основаниях. Его должно решить эксперимент. Мы призываем коллаборации LHC исследовать данные по дилептонам и диджетам с высокими массами на предмет наличия подавления, предсказываемого форм-фактором из ур.~(\ref{eq:F}) с \(\LUV \sim 9\) ТэВ. Нулевой результат исключит сценарий B, и хотя сценарий A остается логически возможным, это ослабит позиции YUCT как решения проблемы иерархии.
	
	\section*{Благодарности}
	Автор благодарит участников семинара YUCT за критические обсуждения.
	
	\section*{Доступность данных}
	Все вычисления являются самодостаточными; код Python, использованный для численных оценок, доступен в репозитории YPSDC: \url{https://github.com/Alexey-Yakushev-YUCT/YPSDC}.
	
	\begin{thebibliography}{99}
		\bibitem{AppendixL} A.\,V.\,Yakushev, ``Приложение L: Фрактальное масштабирование ошибок координации,'' в \emph{YUCT}, Zenodo (2026).
		\bibitem{AppendixV} A.\,V.\,Yakushev, ``Приложение V: Время как энтропия координационных процессов,'' в \emph{YUCT}, Zenodo (2026).
		\bibitem{AppendixG} A.\,V.\,Yakushev, ``Приложение G: Гравитация как геометрическое натяжение координационной сети и циклическая космология в YUCT,'' в \emph{YUCT}, Zenodo (2026).
		\bibitem{AppendixM} A.\,V.\,Yakushev, ``Приложение M: Космология YUCT — инфляция как координационный фазовый переход,'' в \emph{YUCT}, Zenodo (2026).
		\bibitem{AppendixLambda} A.\,V.\,Yakushev, ``Приложение \(\Lambda\): Разрешение проблем иерархии и космологической постоянной в YUCT,'' в \emph{YUCT}, Zenodo (2026).
		\bibitem{AppendixJ} A.\,V.\,Yakushev, ``Приложение J: Полный вывод параметров вкуса Стандартной модели из топологических смещений YUCT,'' в \emph{YUCT}, Zenodo (2026).
		\bibitem{AppendixFerra} A.\,V.\,Yakushev, ``Приложение Ferra1.2: Универсальные координационные константы для упорядоченных и неупорядоченных фаз,'' в \emph{YUCT}, Zenodo (2026).
	\end{thebibliography}
	
\end{document}